\scanpage{163}
% archive sheet number: 162
% L1-2
\[
\underline{\operatorname{Hom}}(G,H)=0\quad\text{si}\quad
\left\{\begin{array}{l}G\text{ à fibres géom connexes}\\H\text{ non ramifié sur }S\end{array}\right.
\]
% L3
\[
\underline{\operatorname{Hom}}(\mathbb Z/n\mathbb Z,G)\simeq {}_nG\qquad\text{d'où lignes 1,2}
\]
% L4
\[
\underline{\operatorname{Hom}}(\mathbb Z,G)\simeq G\qquad\text{d'où ligne 7}
\]
% L5-8
\diagram{\begin{tikzcd}[ampersand replacement=\&,column sep=small,row sep=small]
\underline{\operatorname{Hom}}(\alpha_p,G)\arrow[d,phantom,"\simeq"{rotate=90}]\&\simeq\operatorname{Ker}\bigl(\underline V(\omega_G)\xrightarrow{\text{puiss. }p.\text{ième}}\underline V(\omega_G)\bigr)\\
\underline{\operatorname{Hom}}(\alpha_p,\underline{\unsure{\mathscr U|\mathcal U}}_G^{(1)})\&\begin{gathered}\underline{\unsure{t|\mathfrak t}}_G\\\text{schéma de Lie de }G\end{gathered}
\end{tikzcd}}
\margin{d'où ligne 3 moins l'avant-dernière ----}
% L9
\[
\underline{\operatorname{Hom}}({\unsure{\mu|H}}_p,G)\simeq\operatorname{Ker}\bigl(\underline V\omega_G)\xrightarrow{\text{id-puiss. }p\text{ième}}\underline V\omega_G)\bigr)\qquad\text{d'où ligne 4}
\]
% L10-11
$\underline{\operatorname{Hom}}({\unsure{\mathbb G|G}}_m,G)$ est un foncteur ``non ramifié'' d'où ligne 5

% L12-13
$\underline{\operatorname{Hom}}(A,G)$ est un foncteur ``non ramifié'' d'où (partiellement) ligne 6.

% L14
\[
\underline{\operatorname{Hom}}(\unsure{\mathbb V|V},G)\simeq\underline V(\omega_G)\qquad\text{d'où ligne 8}
\]

% L15-18
Cas difficiles
\diagram{$\begin{array}{l|l}
a)\quad\underline{\operatorname{Hom}}(\mathbb G_a,G)&b)\quad\underline{\operatorname{Hom}}(G,\mathbb V)\\
G=\alpha_p,\mathbb G_a,\mathbb V,\ //A&G=\alpha_p,\mathbb G_a,\mathbb V_{\unsure{k|a}},\ //A
\end{array}$}
(implique b) pour $G=\mathbb G_a,\mathbb V,\alpha_p$)

% L19-20
Il y a essentiellement \deleted{deux}\inserted{quatre} types ``délicats'' ------
